% CVPR 2026 Paper Template; see https://github.com/cvpr-org/author-kit

\documentclass[10pt,twocolumn,letterpaper]{article}

%%%%%%%%% PAPER TYPE  - PLEASE UPDATE FOR FINAL VERSION
% \usepackage{cvpr}              % To produce the CAMERA-READY version
\usepackage[review]{cvpr}      % To produce the REVIEW version
% \usepackage[pagenumbers]{cvpr} % To force page numbers, e.g. for an arXiv version
%Add more
\usepackage{algorithm}

% Add more
\usepackage{algpseudocode} 
\usepackage{amsmath}
\usepackage{multirow}
\usepackage{amsfonts} % For \mathbb
\usepackage{comment}
%
% These are recommended to typeset listings but not required. See the subsubsection on listing. Remove this block if you don't have listings in your paper.
\usepackage{newfloat}
\usepackage{listings}
% Import additional packages in the preamble file, before hyperref
%% This file contains a number of tweaks that are typically applied to the main document.
%% They are not enabled by default, but can be enabled by uncommenting the relevant lines.

%%
%% Inline annotations; for predefined colors, refer to "dvipsnames" in the xcolor package:
%% https://tinyurl.com/overleaf-colors
%%
\newcommand{\red}[1]{{\color{red}#1}}
\newcommand{\todo}[1]{{\color{red}#1}}
\newcommand{\TODO}[1]{\textbf{\color{red}[TODO: #1]}}
\newcommand{\NMD}[1]{\textbf{\color{blue}[NMD: #1]}}

%%
%% disable for camera ready / submission by uncommenting these lines  
%%
% \renewcommand{\TODO}[1]{}
% \renewcommand{\todo}[1]{#1}

%%
%% work harder in optimizing text layout. Typically shrinks text by 1/6 of page, enable
%% it at the very end of the writing process, when you are just above the page limit
%%
% \usepackage{microtype}

%%
%% fine-tune paragraph spacing
%%
% \renewcommand{\paragraph}[1]{\vspace{.5em}\noindent\textbf{#1.}}

%%
%% globally adjusts space between figure and caption
%%
% \setlength{\abovecaptionskip}{.5em}


%%
%% Allows "the use of \paper to refer to the project name"
%% with automatic management of space at the end of the word
%%
% \usepackage{xspace}
% \newcommand{\paper}{ProjectName\xspace}

%%
%% Commonly used math definitions
%%
% \DeclareMathOperator*{\argmin}{arg\,min}
% \DeclareMathOperator*{\argmax}{arg\,max}

%%
%% Tigthen underline
%%
% \usepackage{soul}
% \setuldepth{foobar}

% It is strongly recommended to use hyperref, especially for the review version.
% hyperref with option pagebackref eases the reviewers' job.
% Please disable hyperref *only* if you encounter grave issues, 
% e.g. with the file validation for the camera-ready version.
%
% If you comment hyperref and then uncomment it, you should delete *.aux before re-running LaTeX.
% (Or just hit 'q' on the first LaTeX run, let it finish, and you should be clear).
\definecolor{cvprblue}{rgb}{0.21,0.49,0.74}
\usepackage[pagebackref,breaklinks,colorlinks,allcolors=cvprblue]{hyperref}

% \newcommand{\TODO}[1]{\textbf{\color{red}[TODO: #1]}}
% \newcommand{\NMD}[1]{\textbf{\color{blue}[NMD: #1]}}
\newcommand{\tb}[1]{\textcolor{blue}{#1}}

%%%%%%%%% PAPER ID  - PLEASE UPDATE
\def\paperID{22685} % *** Enter the Paper ID here
\def\confName{CVPR}
\def\confYear{2026}

%%%%%%%%% TITLE - PLEASE UPDATE
\title{H-SFP: Lightweight Federated Learning via Hierarchical Split Prototyping}

%%%%%%%%% AUTHORS - PLEASE UPDATE
\author{First Author\\
Institution1\\
Institution1 address\\
{\tt\small firstauthor@i1.org}
% For a paper whose authors are all at the same institution,
% omit the following lines up until the closing ``}''.
% Additional authors and addresses can be added with ``\and'',
% just like the second author.
% To save space, use either the email address or home page, not both
\and
Second Author\\
Institution2\\
First line of institution2 address\\
{\tt\small secondauthor@i2.org}
}

\begin{document}
\maketitle

\begin{abstract}
Federated learning (FL) offers privacy-preserving distributed training but struggles in real-world settings where non-IID data, device heterogeneity, and high communication costs jointly degrade performance. While split learning (SL) alleviates local computation by offloading part of the model to the server, it incurs heavy communication overhead from transmitting intermediate activations. Hierarchical FL (HFL) improves scalability across edge-cloud hierarchies but fails to address these communication and statistical challenges.
To overcome these limitations, we propose hierarchical split-federated prototyping (H-SFP),  a unified framework that jointly tackles statistical heterogeneity, device heterogeneity, and communication bottlenecks. H-SFP hierarchically partitions the model across clients, edge servers, and the cloud to minimize local computation, and replaces costly activation or gradient exchanges with compact class prototypes that capture semantic information. These prototypes are aggregated first within edge clusters and then globally, enabling stable, semantically aligned updates that regularize local training and suppress client drift.
We provide convergence guarantees for non-convex objectives and extensive empirical validation on CIFAR and HAM10000 datasets. H-SFP consistently delivers higher accuracy under severe non-IID conditions, while achieving over $10\times$ lower communication cost than hybrid baselines and over $300\times$ lower than conventional FL. Our results demonstrate that prototype-driven hierarchical coordination can unlock scalable, communication-efficient, and statistically robust federated training.
\end{abstract}    
\section{Introduction}
\label{sec:intro} 
    Federated Learning (FL) \cite{pmlr-v54-mcmahan17a} faces a trilemma \cite{9084352,MAL-083}: client drift from non-IID data \cite{9084352, GUPTA20181}, system heterogeneity, and high communication overhead \cite{MAL-083}. While FedProx \cite{9084352, dai2023tacklingFedProx} targets non-IID data, it still requires full-model training, which is infeasible for resource-constrained devices \cite{10825769}. Moreover, split Learning (SL) \cite{GUPTA20181, vepakomma2019split} solves the computational burden but has a sequential bottleneck \cite{thapa2022splitfed}. Split FL (SFL) \cite{thapa2022splitfed} parallelizes SL but introduces scalability bottlenecks \cite{10980018} and prohibitive communication/privacy overhead from exchanging high-dimensional activations \cite{11091500, zhang2025amperecommunicationefficienthighaccuracysplit, qiu2023evaluating,wang2025towards}. We propose H-SFP (Hierarchical Split-Federated Prototyping) to address these limitations. H-SFP integrates: (i) a hierarchical topology \cite{10980018, 9148862} for scalability; (ii) split execution \cite{GUPTA20181} for low client computation; and (iii) prototype-based communication \cite{tan2022fedproto}, replacing costly activation/gradient exchange \cite{11091500} with lightweight and abstract prototypes. 
    
    This integration yields a crucial synergistic benefit. While FedProto \cite{tan2022fedproto} shows surprising robustness to non-IID data (as our results in \cref{tab:main_accuracy_results} confirm), its flat, centralized architecture does not address server-side scalability. Our work builds on this, positioning H-SFP to not only surpass FedProto's state-of-the-art accuracy but also to simultaneously solve the scalability bottleneck. By performing ``intragroup prototype aggregation" at the edge-server level, H-SFP refines and stabilizes representations, filtering client-level noise before global aggregation. This makes the framework uniquely robust to non-IID data \cite{9084352, GUPTA20181}.

    Our primary contributions are: 
    \begin{itemize} 
        \item We propose H-SFP, a novel three-tier architecture that synergistically integrates split learning, a hierarchical topology, and prototype-based communication to to solve distributed learning bottlenecks holistically. 
        \item We prove hierarchical aggregation acts as a stabilization mechanism, making ours robust to the high statistical heterogeneity challenging standard prototype methods. 
        \item We provide extensive empirical validation of H-SFP's efficiency and scalability against standard FL, HierFL, SFL, HSFL, contrastive-based FL, and FedProto baselines. 
    \end{itemize}
\section{Background and Related Works}
    \begin{figure}
        \centering
        \includegraphics[width=\linewidth]{CVPR-2026/Figures/Related Works.pdf}
        \caption{Overview of FL paradigms. (A) Conventional FL exchanges full models with the cloud. (B) SplitFL transmits smashed data but depends on cloud gradients. (C) Hierarchical FL uses edge servers for intermediate aggregation. (D) Hierarchical SplitFL combines model splitting with hierarchical aggregation.}
        \label{fig:placeholder}
    \end{figure}
    Our work is positioned at the intersection of three key research areas aimed at solving the FL trilemma. The first, statistical heterogeneity, causes ``client drift" in FedAvg \cite{pmlr-v54-mcmahan17a}. Major solutions include regularization methods like FedProx \cite{dai2023tacklingFedProx}, drift-correction via control variates as in SCAFFOLD \cite{pmlr-v119-karimireddy20a}, and dynamic regularizers like FedDyn \cite{durmus2021federated}. A different approach, personalized FL (PFL), embraces heterogeneity. This is often achieved by decoupling models, as seen in FedPer \cite{arivazhagan2019federatedlearningpersonalizationlayers} and FedRep \cite{collins2021exploiting}, which train a shared representation but private heads. Other PFL strategies include pFedMe \cite{t2020personalized} using Moreau envelopes, and MOON \cite{9578660} using contrastive learning to align representations.

The second area, system heterogeneity, addresses client-side constraints. SL \cite{GUPTA20181, vepakomma2019split} solves this by partitioning the model, forcing clients to train only lightweight initial layers. To solve SL's sequential bottleneck, SFL \cite{thapa2022splitfed} was proposed to parallelize clients. A major alternative for handling computational diversity is HeteroFL \cite{diao2021heterofl}, which trains different-sized sub-models. System heterogeneity in terms of clients performing variable local work is addressed by methods like FedNova \cite{10.5555/3495724.3496362}, which normalizes local updates.

The third area tackles communication and scalability bottlenecks. Communication overhead is often reduced by model compression, such as quantization in QSGD \cite{alistarh2017qsgd} or sparsification in Deep Gradient Compression \cite{lin2017deep}. Another strategy is Federated Knowledge Distillation (e.g., FedDF \cite{lin2020ensemble}, FedGen \cite{venkateswaran2023fedgen}). The strategy we adopt, Prototype-Based FL \cite{tan2022fedproto}, exchanges lightweight prototypes, which both saves bandwidth and aids personalization. Separately, the single-server scalability bottleneck is solved by Hierarchical FL (HFL) \cite{rana2022hierarchical,9148862,9488756}, which introduces edge servers for intermediate aggregation and cloud workload reduction.

The existing literature has largely explored these techniques as separate solutions, each with their own limitations. For instance, SFL creates new server bottlenecks \cite{10980018} and incurs high communication costs from exchanging high-dimensional activations \cite{11091500, zhang2025amperecommunicationefficienthighaccuracysplit}, while FedProto can suffer from prototype drift under high heterogeneity \cite{bensiah2025feddpa, zhang2020personalized}. To our knowledge, no prior work has proposed a unified framework that synergistically integrates all three. Our H-SFP framework fills this gap. We use the HFL structure to solve SFL's scalability bottleneck and replace SFL's costly activation exchange with lightweight prototypes. 
It is critical to distinguish ours from methods like FedProto \cite{tan2022fedproto}. While FedProto uses prototypes as regularizers to align local training, H-SFP's prototypes are generative statistics ($\mu, \sigma$) used to synthesize new training data for the upper tiers, thereby enabling a fully decoupled, feed-forward training process. This yields a crucial synergistic benefit: the hierarchical, intragroup aggregation stabilizes the prototypes, mitigating drift. Therefore, our framework is simultaneously robust to non-IID data, computationally efficient for clients, and scalable to massive networks.


\section{The Proposed H-SFP Framework}
\label{sec:methodology}
We propose H-SFP, a three-tier framework designed to holistically address the FL trilemma: statistical heterogeneity, system heterogeneity, and communication bottlenecks. H-SFP's core innovation is a decoupled, generative workflow that severs the tight coupling of end-to-end backpropagation. Each tier (client, edge, cloud) trains independently: it ``packages'' its learned representations into lightweight feature statistics (mean and standard deviation) and transmits only these statistics upward. Additionally, the concrete algorithms for each tier are presented in Appendix A.

\begin{figure*}[t]
    \centering
    \includegraphics[width=\linewidth]{CVPR-2026/Figures/method-2.pdf}
    \caption{The H-SFP Hierarchical Architecture and Data Flow. Our framework is a three-tier system (Client, Edge, Cloud). Clients perform local training (detailed in \cref{fig:method-p1}) and ``package" their feature statistics to be sent to an Edge Server. The Edge Server uses these statistics to synthesize its own training data, trains its model segment, and then packages its own feature statistics for the Cloud. The Cloud repeats the synthesis process to train the final, task-specific model layers, eliminating the need for end-to-end backpropagation.}
    \label{fig:method-p2}
\end{figure*}

% \TODO{From my oppinion, 3.2 is a little bit verbose, 3.4 is redundant.}

\subsection{System Architecture and Model Partitioning}
\label{subsec:system_model}
% \TODO{In this part, please focus on Problem definition and system architecture. Please follow 2.1 of https://arxiv.org/pdf/2409.18448 
% Remember to note the three-tier or something else of yours, with your notation definitions.
% }

\begin{figure}[t]
    \centering
    \includegraphics[width=\linewidth]{CVPR-2026/Figures/method-3.pdf}
    \caption{Model Partitioning in H-SFP. The global model $W$ is strategically split into three disjoint segments corresponding to the architectural tiers: the lightweight client-side model ($W_c$), the intermediate edge-side model ($W_e$), and the final cloud-side model ($W_g$). This partitioning offloads the primary computational burden from resource-constrained clients.}
    \label{fig:method-p3}
\end{figure}
We consider a three-tier ``cloud-edge-client'' architecture (\cref{fig:method-p2}), comprising a global cloud server, $M$ edge servers, and $K$ clients. Clients are partitioned into $M$ disjoint clusters $\{\mathcal{C}_m\}_{m=1}^M$, with each client $k \in \mathcal{C}_m$ associated with its parent edge server $m$. This hierarchical topology directly addresses the scalability and single-point-of-failure bottlenecks of centralized FL. Instead of a single cloud bottleneck processing updates from all $N$ clients, each of the $M$ edge servers performs intermediate aggregation for its local cluster. The cloud, in turn, only aggregates the $M$ processed results from the edges, significantly reducing communication overhead and latency. To address system heterogeneity (i.e., client computational constraints), H-SFP integrates split learning \cite{GUPTA20181, vepakomma2019split}. As shown in \cref{fig:method-p3}, the global model $W$ is partitioned into three disjoint segments corresponding to the tiers: lightweight client-side extractors ($W_c$), intermediate edge-side models ($W_e$), and final cloud-side layers ($W_g$). This partitioning strategy solves the computational bottleneck, as resource-constrained clients are only responsible for executing the small $W_c$ segment.

\subsection{Motivation}
    \begin{figure}[h]
        \centering
        \includegraphics[width=\linewidth]{CVPR-2026/Figures/method-1.pdf}
        \caption{The Core Generative Prototyping Mechanism. This diagram illustrates the process at each tier. (Left) On clients, self-supervised contrastive learning (pulling similar, pushing dissimilar) is used to learn robust feature representations. (Right) Instead of communicating high-dimensional gradients, clients ``package" these representations into lightweight statistical moments (mean $\mu$ and std $\sigma$). These statistics (prototypes) are then used by the next tier to ``synthesize" new feature data, enabling decoupled and communication-efficient training.}
        \label{fig:method-p1}
    \end{figure}
    While SL successfully addresses the computational bottleneck, it introduces a critical communication bottleneck. In SFL, the model segments remain tightly coupled, requiring a full end-to-end backpropagation pass for every iteration. This necessitates the frequent, high-volume transmission of intermediate activations and gradients, a process that is impractical and defeats the goal of an efficient FL. To make communication efficient, two key challenges must be addressed: (i) reducing the training of the subnetworks to remove the synchronous, end-to-end backpropagation, and (ii) an alternative form of knowledge transfer that is lightweight yet semantically meaningful.
    
    Therefore, we propose a novel multi-tier generative prototyping mechanism, as illustrated in \cref{fig:method-p2} and \cref{fig:method-p1}. This feed-forward workflow completely decouples the training of the model segments ($W_c, W_e, W_g$), solving the communication and synchronization bottlenecks of SFL. To prevent long-term client drift, which can degrade purely generative approaches, our framework integrates this fast-timescale prototype exchange with a slow-timescale periodic averaging of the client-side feature extractors ($W_c$). This creates a robust hybrid system that is both communication-efficient and stable. In \cref{sec:multi_tier_learning}, we detail the independent, self-supervised learning objective for each tier. In \cref{sec:transmission_augmentation}, we describe the core mechanics of this workflow: how each tier ``packages'' its learned representations into lightweight statistics, and how the next tier ``synthesizes" a new feature dataset, $D_{syn}$, from these statistics to serve as its own training input.

\subsection{Multi-tier Self-Learning} 
\label{sec:multi_tier_learning}
    H-SFP decouples the training of the network segments, enabling each tier to learn independently. This is achieved via a multi-tier generative process where feature statistics, not gradients, are passed unidirectionally up the hierarchy.
    
    The first tier of learning occurs on the resource-constrained clients. As shown in \cref{fig:method-p1}, each client $k$ trains its local feature extractor $W_c$ using only its private data $\mathcal{D}_k$. The objective $\mathcal{L}_k(W_c)$ is to learn a high-quality, robust representation space using a self-supervised contrastive loss. Specifically, we use the NT-Xent loss \cite{chen2020simple}. For a batch $B$, let $z_i = W_c(x_i)$ be the feature embedding for a sample $x_i$:
    \begin{multline}
        \label{eq:supcon}
        \mathcal{L}_k(W_c) = \mathcal{L}_{\text{con}, k} \\
        = \sum_{i \in B} \frac{-1}{|P(i)|} \sum_{p \in P(i)} \log \frac{\exp(z_i \cdot z_p / \tau)}{\sum_{a \in A(i)} \exp(z_i \cdot z_a / \tau)},
    \end{multline}
    where $A(i) = B \setminus \{i\}$ is the set of all other samples in the batch, $P(i) \equiv \{p \in A(i) | y_p = y_i\}$ is the set of positive samples in the batch (i.e., those with the same class label as $i$), and $\tau$ is a temperature hyperparameter. This encourages the client's model to learn a feature space where samples from the same class are pulled together, while samples from different classes are pushed apart. After local training, the client distills its learned feature space into lightweight class-wise feature statistics (e.g., class-conditional mean $\mu_k$ and standard deviation $\sigma_k$). These statistics are the only information communicated upward to the edge server (in \cref{fig:method-p2}).
    
    Then, the second tier (edge server $m$) operates without access to any real data, enabling a fully decoupled architecture. First, the edge server aggregates the statistics $\{\mu_k, \sigma_k\}$ received from its client cluster $\mathcal{C}_m$ to form a stable, cluster-level feature distribution $(\mu_m, \sigma_m)$. Second, the subsequent training process is nearly identical to the client's, but it operates on synthetically generated input. The edge server unpacks these aggregated statistics and synthesizes a new feature dataset, $D_{syn}^{edge}$. This synthetic data, generated by sampling from the cluster's distribution (using $\mu_m$ as the signal and $\sigma_m$ for perturbation), serves as the input for the edge-side model, $W_e$. The edge server then trains $W_e$ by applying the same contrastive loss (defined in \cref{eq:supcon}) on its own output representations:
    \begin{equation}
        \mathcal{L}_m(W_e) = \mathcal{L}_{\text{con}, m}(W_e(D_{syn}^{edge})).
    \end{equation}
    We deliberately select this contrastive objective over a self-supervised loss to enforce objective consistency across tiers and to learn the relative relationships between feature clusters rather than overfitting the synthetic Gaussian distributions. This results in a more robust and generalized representation for the final cloud-tier classifier. This step learns an effective transformation of the client-level features into a new, more refined representation space. Like the clients, the edge server then packages the statistics of its output features and forwards them to the cloud.
    
    Finally, the final tier repeats the synthesis process, aggregating the feature statistics from all $\mathcal{M}$ edge servers to create a global distribution and generate its own synthetic dataset, $D_{syn}^{global}$. However, its objective is different. As shown in ~\cref{fig:method-p2}, the cloud server trains the final task-specific model $W_g$ (e.g., the classifier head) using a standard supervised loss, such as cross-entropy:
    \begin{equation}
        \mathcal{L}_g(W_g) = \mathcal{L}_{\text{CE}}(W_g(D_{syn}^{global})).
    \end{equation}
    This multi-tier, generative, and feed-forward process completely breaks the reliance on end-to-end backpropagation, solving the limitations of SL-based approaches.
    

\subsection{Prototype Transmission and Augmentation}
\label{sec:transmission_augmentation}
    % --- REVISED ---
    H-SFP's novelty lies in its communication, which replaces backpropagation with a feed-forward, generative process. This involves two stages: transmitting lightweight statistics and augmenting data at the receiving tier.
    % --- END REVISED ---

    \begin{algorithm}[!h]
        \caption{H-SFP: Client-Side Training}
        \label{alg:client}
        \begin{algorithmic}[1]
        \State \textbf{Input:} Client $k$, local data $D_k$, learning rate $\eta$, epochs $E$.
        \State \textbf{Receive:} Client model $W_c$.
        \State Set $W_{c,k} \leftarrow W_c$.
        \For{epoch $e = 1$ to $E$}
            \For{batch $(x, y)$ in $D_k$}
                \State $z \leftarrow W_{c,k}(x)$ \Comment{Compute feature embeddings}
                \State $\mathcal{L}_{\text{con}} \leftarrow \text{ContrastiveLoss}(z, y)$ \Comment{See \cref{eq:supcon}}
                % <<< REVISION: Removed L_align steps
                \State $W_{c,k} \leftarrow W_{c,k} - \eta \nabla \mathcal{L}_{\text{con}}$
            \EndFor
        \EndFor
        \State Compute final local statistics $\text{Stats}_k = \{(\mu_k^{(j)}, \sigma_k^{(j)}) \mid j \in \mathcal{Y}_k\}$
        \Comment{See \cref{mean_client,std_client}}
        \State \textbf{Return:} $W_{c,k}$, $\text{Stats}_k$ to Edge Server $m$.
        \end{algorithmic}
    \end{algorithm}
    
    \subsubsection{Prototype Transmission}
    % --- REVISED ---
        Our communication mechanism is a dual-timescale hybrid, which explains the distinct cost categories in \cref{tab:comm_overhead}. Most rounds (fast timescale) involve prototype exchange, where clients send only lightweight statistics ($\mu, \sigma$) instead of costly activations. This replaces the standard SFL bottleneck and corresponds to the large `Client smashed upload' cost. Infrequently (slow timescale, interval $I_c$), model averaging of the client parameters ($W_{c,k}$) occurs in the background to combat long-term drift, corresponding to the tiny 'Client model upload' cost. A similar slow-timescale process occurs at the interval $I_e$ between the edge and cloud.

    % --- END REVISED ---
        \begin{equation}
            \mu_k^{(j)} = \frac{1}{|\mathcal{D}_k^{(j)}|} \sum_{x_i \in \mathcal{D}_k^{(j)}} W_{c,k}(x_i),
            \label{mean_client}
        \end{equation}
        \begin{equation}
            \sigma_k^{(j)} = \sqrt{\frac{1}{|\mathcal{D}_k^{(j)}|} \sum_{x_i \in \mathcal{D}_k^{(j)}} (W_{c,k}(x_i) - \mu_k^{(j)})^2}.
            \label{std_client}
        \end{equation}
    % --- REVISED ---
        The transmitted package $(\mu_k^{(j)}, \sigma_k^{(j)})$ represents the complete generative statistics for class $j$. This mechanism drastically reduces communication overhead. The cost per client is $O(J_k \times d_c)$, where $J_k$ is the class count and $d_c$ is the feature dimension. This cost is independent of dataset size and model parameters $P_c$, contrasting sharply with FedAvg's $O(P_c)$ cost and standard SFL's $O(d_c)$ cost per batch.
    % --- END REVISED ---

    \subsubsection{Prototype Augmentation}
        \begin{figure}[h]
            \centering
            % Use the path to the PDF you generated
            \includegraphics[width=\linewidth]{CVPR-2026/Figures/synthesis_comparison.pdf}
            \caption{Visualizing the necessity of dimension-aware generative augmentation. Real-world feature distributions are not uniform hyperballs, as variance differs across feature dimensions. This figure illustrates the impact of this fact. \textbf{(a) Naive Synthesis:} A naive approach, using only the reference vector (mean $\mu$) and simple, spherical noise, fails to capture the true, complex shape of the data distribution. \textbf{(b) Our Method:} ``Unpacking" the full transmitted statistics (both the mean $\mu$ and the standard deviation vector $\sigma$) allows for dimension-aware perturbation, which correctly reconstructs the true, elliptical shape of the feature clusters.}
            \label{fig:synthesis_comparison}
        \end{figure}
    
    % --- REVISED ---
        Upon receiving statistics $\{(\mu_k, \sigma_k)\}_{k \in \mathcal{C}_m}$, the edge server $m$ aggregates them into a cluster-level distribution:
        \begin{equation}
            \mu_m^{(j)} = \frac{1}{|\mathcal{C}_m|} \sum_{k \in \mathcal{C}_m} \mu_k^{(j)},
        \end{equation}
        \begin{equation}
            (\sigma_m^{(j)})^2 = \frac{1}{|\mathcal{C}_m|} \sum_{k \in \mathcal{C}_m} (\sigma_k^{(j)})^2.
        \end{equation}
        This distribution acts as a generative model to synthesize training data $D_{syn}^{edge}$ for the server's model segment $W_e$. As visualized in \cref{fig:synthesis_comparison}, a naive synthesis using only the mean $\mu_m^{(j)}$ and spherical noise fails, as it ignores the non-uniform, elliptical shape of true feature distributions. Our method, inspired by prototype augmentation \cite{zhu2021prototype}, recognizes this. It ``unpacks" the full statistics, using the standard deviation vector $\sigma_m^{(j)}$ to perform dimension-aware ``prototype perturbation". Specifically, a synthetic feature $z_{syn}$ for class $j$ is generated by sampling from a multivariate Gaussian by
        \begin{equation}
            \label{eq:synthesis}
            z_{syn} \sim \mathcal{N}(\mu_m^{(j)}, \text{diag}(\sigma_m^{(j)} + \epsilon)^2),
        \end{equation}
        where $\text{diag}(\cdot)$ creates a diagonal covariance matrix to model the elliptical clusters, and $\epsilon$ is a small stability constant. We intentionally use a diagonal covariance matrix as a trade-off between modeling capacity and communication cost. Our ablation study (\cref{tab:ablation}, F vs. D) confirms that modeling per-dimension variance (transmitting $\sigma$) is essential. However, transmitting a full covariance matrix would increase the communication cost from $O(d_c)$ to $O(d_c^2)$, re-introducing a major bottleneck. Our approach captures the elliptical feature distributions (\cref{fig:synthesis_comparison}) without this prohibitive cost.

    \subsection{Convergence Analysis}
        We provide convergence guarantees for H-SFP, with the full proof detailed in Appendix C. Our analysis differs from standard FL as it must account for our decoupled, generative workflow. We adopt a hierarchical, bottom-up strategy under standard non-convex assumptions (e.g., L-smoothness) and a novel $L_{stat}$-Lipschitz assumption to bound the impact of changing generative statistics. Moreover, the core of the proof is to show that the upper-tier models (edge servers and cloud) converge even while training on a ``moving target" synthetic dataset. We first prove that the client-tier statistics ($\mu_k, \sigma_k$) stabilize as the client models ($W_c$) converge. This stabilization ensures the ``distribution shock" at the edge servers diminishes ($\delta_m^t \rightarrow 0$), allowing the edge model ($W_e$) to converge. This logic extends to the cloud, confirming that the average squared gradient norm at each tier converges to a bounded neighborhood, thus proving the stability of our framework.
\begin{table*}
    \centering
    \caption{F1-Score (\%) on CIFAR-10, CIFAR-100, and HAM10000 under IID and non-IID (Dirichlet) settings. (A), (B), and (C) represent client-edge ($I_c$) and edge-cloud ($I_e$) aggregation intervals of (5, 10), (10, 20), and (25, 50), respectively. \textbf{Bold} indicates the best performance, and \underline{underline} indicates the second-best.}
    \begin{tabular}{l|ccc|ccc|ccc}
        \hline
        \multirow{2}{*}{\textbf{Method}} &
        \multicolumn{3}{c|}{\textbf{CIFAR-10}} &
        \multicolumn{3}{c|}{\textbf{CIFAR-100}} &
        \multicolumn{3}{c}{\textbf{HAM10000}} \\
        \cline{2-10}
        & \textbf{IID} & $dir(0.3)$ & $dir(0.7)$ & \textbf{IID} & $dir(0.3)$ & $dir(0.7)$ & \textbf{IID} & $dir(0.3)$ & $dir(0.7)$ \\
        \hline
        FedAvg & 57.25\% & 12.15\% & 23.60\% & 48.85\% & 2.10\% & 5.40\% & 68.28\% & 9.10\% & 13.45\%\\
        FedSGD & 18.80\% & 5.80\% & 6.40\% & 10.75\% & 1.80\% & 4.35\% & 25.85\% & 4.15\% & 6.34\%\\
        FedNova & 23.50\% & 14.20\% & 18.35\% & 50.24\% & 3.15\% & 6.80\% & 69.60\% & 10.05\% & 14.35\%\\
        FedProx & 58.25\% & 16.50\% & 25.40\% & 52.75\% & 4.05\% & 7.20\% & 71.15\% & 11.20\% & 15.45\%\\
        \hline
        FedProto & 55.10\% & \underline{38.50\%} & 45.20\% & 46.50\% & 6.50\% & 8.80\% & 67.20\% & 28.10\% & 35.00\%\\
        \hline
        HierFL & 62.50\% & 13.10\% & 23.60\% & 30.15\% & 1.95\% & 4.78\% & 60.15\% & 16.40\% & 20.85\%\\
        \hline
        SplitFed & 61.85\% & 31.20\% & 39.95\% & 35.00\% & 4.15\% & 7.05\% & 65.78\% & 20.25\% & 25.15\%\\
        \hline
        FedGen & 59.80\% & 37.90\% & 46.10\% & 51.20\% & 6.70\% & 9.15\% & 72.50\% & 29.20\% & 35.80\%\\
        FedDF & 61.50\% & 32.00\% & 40.50\% & 53.10\% & 4.50\% & 7.80\% & 73.10\% & 22.50\% & 27.80\% \\
        \hline
        HSFL (A) & 51.35\% & 17.80\% & 26.54\% & 35.47\% & 5.10\% & 8.25\% & 68.82\% & 25.30\% & 30.45\%\\
        HSFL (B) & 50.80\% & 17.15\% & 25.90\% & 33.58\% & 4.80\% & 7.65\% & 67.74\% & 24.10\% & 28.83\%\\
        HSFL (C) & 49.18\% & 16.90\% & 25.40\% & 32.65\% & 4.10\% & 6.95\% & 66.98\% & 23.05\% & 27.95\%\\
        \hline
        \textbf{Ours (A)} & \textbf{67.44\%} & \textbf{40.15\%} & \textbf{48.85\%} & \textbf{55.10\%} & \textbf{7.15\%} & \textbf{10.48\%} & \textbf{80.86\%} & \textbf{31.80\%} & \textbf{37.50\%}\\
        \textbf{Ours (B)} & \underline{66.80\%} & 39.80\% & \underline{48.15\%} & \underline{54.25\%} & \underline{6.80\%} & \underline{9.95\%} & \underline{79.96\%} & \underline{31.10\%} & \underline{36.90\%}\\
        \textbf{Ours (C)} & 65.44\% & 39.10\% & 47.50\% & 53.65\% & 6.15\% & 9.05\% & 79.15\% & 30.25\% & 36.15\%\\
        \hline
    \end{tabular}
    \label{tab:main_accuracy_results}
\end{table*}
\section{Experiments}
\label{sec:experiments}
    We evaluate H-SFP on CIFAR-10/100 and HAM10000 against centralized (FedAvg, FedProx), prototype (FedProto), and hybrid (HierFL, SplitFed, HSFL) baselines. HSFL serves as our primary architectural baseline. We test IID and non-IID (Dirichlet $dir(0.3)$/$dir(0.7)$) settings with 200 clients grouped under 10 edge servers for hierarchical methods. Models use an Adam optimizer ($lr=1e-4$), and we report F1 score, communication cost, and client RAM. Additionally, the experiments are decribed in Appendix B.
\subsection{Main Results: Model Performance}
    \cref{tab:main_accuracy_results} presents the primary classification results. Our proposed H-SFP (A) consistently outperforms all baselines, achieving the highest F1-scores in IID settings (e.g., 80.86\% on HAM10000). The framework's superiority is most evident in the challenging non-IID $dir(0.3)$ scenarios, where standard methods like FedAvg (9.10\%) suffer a catastrophic performance collapse. Notably, H-SFP (A) significantly surpasses not only its closest architectural competitor, HSFL (A) (25.30\% on HAM10000), but also robust, state-of-the-art methods including FedGen (29.20\%) and FedProto (28.10\%). This trend is consistent across CIFAR-10 ($dir(0.3)$), where H-SFP (A) (40.15\%) again leads all methods, including FedProto (38.50\%) and FedGen (37.90\%). This validates that our synergistic design of hierarchical aggregation and generative prototype communication is critically effective at mitigating client drift. Finally, the H-SFP variants (A, B, C) show a graceful trade-off between aggregation frequency and accuracy, with the most frequent (A) performing best.


\subsection{Analysis of Communication Overhead}
\begin{table*}[!ht]
\centering
\caption{Communication overhead (GB) on CIFAR-100 with ResNet-18 over 200 rounds under 200 clients.}
\resizebox{\textwidth}{!}{%
    \begin{tabular}{l|cccccccc|c}
    \hline
    \textbf{Method} &
    \begin{tabular}[c]{@{}c@{}}\textbf{Client}\\ \textbf{smashed}\\ \textbf{upload}\end{tabular} &
    \begin{tabular}[c]{@{}c@{}}\textbf{Edge}\\ \textbf{smashed}\\ \textbf{upload}\end{tabular} &
    \begin{tabular}[c]{@{}c@{}}\textbf{Client}\\ \textbf{model}\\ \textbf{upload}\end{tabular} &
    \begin{tabular}[c]{@{}c@{}}\textbf{Client}\\ \textbf{model}\\ \textbf{download}\end{tabular} &
    \begin{tabular}[c]{@{}c@{}}\textbf{Edge}\\ \textbf{model}\\ \textbf{upload}\end{tabular} &
    \begin{tabular}[c]{@{}c@{}}\textbf{Edge}\\ \textbf{model}\\ \textbf{download}\end{tabular} &
    \begin{tabular}[c]{@{}c@{}}\textbf{Gradient}\\ \textbf{from}\\ \textbf{cloud}\end{tabular} &
    \begin{tabular}[c]{@{}c@{}}\textbf{Gradient}\\ \textbf{from}\\ \textbf{edge}\end{tabular} &
    \textbf{Total (GB)} \\
    \hline
    FedAvg & -- & -- & 351.40 & 3,513.95 & -- & -- & -- & -- & 3,865.35 \\ \hline
    % FedSGD & -- & -- & 351.40 & 3,513.95 & -- & -- & -- & -- & 3,865.35 \\
    % FedNova & -- & -- & 351.40 & 3,513.95 & -- & -- & -- & -- & 3,865.35 \\
    % FedProx & -- & -- & 351.40 & 3,513.95 & -- & -- & -- & -- & 3,865.35 \\ \hline
    HierFL & -- & -- & 351.40 & 3,865.34 & 175.70 & -- & -- & -- & 4,392.44 \\ \hline
    SplitFed & 6.87 & -- & 1.09 & 216.82 & -- & -- & 3.73 & -- & 228.51 \\ \hline
    % --- FEDPROTO ROW ADDED ---
    FedProto & -- & -- & 1.57 & 15.68 & -- & -- & -- & -- & 17.25 \\ \hline
    HSFL ($I_c=5, I_e=10$) & 54.93 & 13.74 & 0.22 & 0.93 & 1.34 & 1.34 & 12.22 & 31.25 & 115.97 \\
    HSFL ($I_c=10, I_e=20$) & 54.93 & 13.74 & 0.11 & 0.48 & 0.67 & 0.67 & 12.22 & 31.25 & 114.07 \\
    HSFL ($I_c=25, I_e=50$) & 54.93 & 13.74 & 0.06 & 0.24 & 0.34 & 0.34 & 12.22 & 31.25 & 113.12 \\ \hline
    \textbf{Ours ($I_c=5, I_e=10$)} & \textbf{9.77} & \textbf{0.19} & \textbf{0.22} & \textbf{0.93} & \textbf{1.34} & \textbf{1.34} & -- & \textbf{--} & \textbf{13.79} \\
    \textbf{Ours ($I_c=10, I_e=20$)} & \textbf{9.77} & \textbf{0.19} & \textbf{0.11} & \textbf{0.48} & \textbf{0.67} & \textbf{0.67} & -- & \textbf{--} & \textbf{11.89} \\
    \textbf{Ours ($I_c=25, I_e=50$)} & \textbf{9.77} & \textbf{0.19} & \textbf{0.06} & \textbf{0.24} & \textbf{0.34} & \textbf{0.34} & -- & \textbf{--} & \textbf{10.94} \\
    \hline
    \end{tabular}
}
\label{tab:comm_overhead}
\end{table*}

\begin{table*}
    \centering
    \caption{F1-Score (\%) on CIFAR-10, CIFAR-100, and HAM10000 under IID scenario. This table shows robustness to data sparsity by varying the total number of clients in a fixed-data experiment. (A), (B), and (C) represent client-edge ($I_c$) and edge-cloud ($I_e$) aggregation intervals of (5, 10), (10, 20), and (25, 50), respectively. \textbf{Bold} indicates the best performance, and \underline{underline} indicates the second-best.}
    \resizebox{\textwidth}{!}{%
        \begin{tabular}{l|cccc|cccc|cccc}
            \hline
            \multirow{2}{*}{\textbf{Method}} &
            \multicolumn{4}{c|}{\textbf{CIFAR-10}} &
            \multicolumn{4}{c|}{\textbf{CIFAR-100}} &
            \multicolumn{4}{c}{\textbf{HAM10000}} \\
            \cline{2-13}
            & \textbf{200} & \textbf{100} & \textbf{50} & \textbf{20} & \textbf{200} & \textbf{100} & \textbf{50} & \textbf{20} & \textbf{200} & \textbf{100} & \textbf{50} & \textbf{20} \\
            \hline
            FedAvg & 57.25\% & 63.10\% & 68.45\% & 72.30\% & 48.85\% & 53.20\% & 57.10\% & 60.50\% & 68.28\% & 73.10\% & 77.05\% & 79.80\%\\
            FedSGD & 18.80\% & 22.50\% & 25.90\% & 28.15\% & 10.75\% & 13.10\% & 15.80\% & 17.90\% & 25.85\% & 29.40\% & 32.15\% & 34.00\%\\
            FedNova & 23.50\% & 28.15\% & 32.40\% & 35.80\% & 50.24\% & 54.80\% & 58.90\% & 61.70\% & 69.60\% & 74.20\% & 78.10\% & 80.60\%\\
            FedProx & 58.25\% & 64.00\% & 69.10\% & 73.00\% & 52.75\% & 57.00\% & 60.90\% & 63.80\% & 71.15\% & 75.90\% & 79.40\% & 81.70\%\\
            \hline
            FedProto & 55.10\% & 62.50\% & 68.10\% & 71.20\% & 46.50\% & 51.90\% & 56.00\% & 59.10\% & 67.20\% & 74.00\% & 78.50\% & 81.10\%\\
            \hline
            HierFL & 62.50\% & 67.10\% & 71.30\% & 74.50\% & 30.15\% & 34.00\% & 37.80\% & 40.10\% & 60.15\% & 65.20\% & 69.80\% & 72.30\%\\
            \hline
            SplitFed & 61.85\% & 66.90\% & 71.05\% & 74.20\% & 35.00\% & 39.10\% & 42.80\% & 45.40\% & 65.78\% & 70.50\% & 74.60\% & 77.90\%\\
            \hline
            FedGen & 59.80\% & 65.00\% & 70.00\% & 73.80\% & 51.20\% & 55.45\% & 59.35\% & 62.25\% & 72.50\% & 77.25\% & 80.75\% & 83.05\% \\
            FedDF & 61.50\% & 66.70\% & 70.90\% & 74.00\% & 53.10\% & 57.35\% & 61.25\% & 64.15\% & 73.10\% & 77.85\% & 81.35\% & 83.65\%\\
            \hline
            HSFL (A) & 51.35\% & 57.10\% & 62.30\% & 66.40\% & 35.47\% & 40.20\% & 44.10\% & 47.00\% & 68.82\% & 74.10\% & 78.20\% & 80.70\%\\
            HSFL (B) & 50.80\% & 56.50\% & 61.90\% & 65.80\% & 33.58\% & 38.10\% & 42.00\% & 44.90\% & 67.74\% & 73.20\% & 77.50\% & 80.10\%\\
            HSFL (C) & 49.18\% & 55.00\% & 60.10\% & 64.20\% & 32.65\% & 37.00\% & 40.80\% & 43.50\% & 66.98\% & 72.50\% & 76.90\% & 79.50\%\\
            \hline
            \textbf{Ours (A)} & \textbf{67.44\%} & \textbf{74.15\%} & \textbf{79.85\%} & \textbf{82.50\%} & \textbf{55.10\%} & \textbf{60.20\%} & \textbf{64.50\%} & \textbf{67.80\%} & \textbf{80.86\%} & \textbf{84.50\%} & \textbf{87.10\%} & \textbf{89.00\%}\\
            \textbf{Ours (B)} & \underline{66.80\%} & \underline{73.80\%} & \underline{79.15\%} & \underline{81.90\%} & \underline{54.25\%} & \underline{59.40\%} & \underline{63.80\%} & \underline{66.90\%} & \underline{79.96\%} & \underline{83.90\%} & \underline{86.70\%} & \underline{88.60\%}\\
            \textbf{Ours (C)} & 65.44\% & 72.50\% & 78.00\% & 80.70\% & 53.65\% & 58.70\% & 62.90\% & 65.90\% & 79.15\% & 83.10\% & 85.90\% & 87.90\%\\
            \hline
        \end{tabular}
    }
    \label{tab:main_robustness_sparsity} 
\end{table*}

\cref{tab:comm_overhead} provides a detailed breakdown of the communication overhead, quantifying the primary advantage of H-SFP. As shown, traditional FL methods like FedAvg are prohibitively expensive ($\sim$3,865 GB) due to full model transfers, and HierFL is even worse (4,392 GB). While hybrid methods like HSFL (113.12 GB) and prototype-based FedProto (17.25 GB) offer significant savings, our H-SFP framework provides the most dramatic reduction.Our most efficient configuration (Ours (C)) requires only 10.94 GB. This is 353x more efficient than FedAvg, 10.4x more efficient than the HSFL baseline, and 1.58x more efficient than the strong FedProto baseline. This gain is achieved by our novel architecture: H-SFP completely eliminates gradient backpropagation (a 43.47 GB cost for HSFL) and replaces the transfer of raw activations (54.93 GB for HSFL) with lightweight feature statistics (only 9.77 GB), validating its ability to resolve the communication bottleneck.

\subsection{Analysis of Client-Side Memory Cost}
\begin{figure}[t]
    \centering
    \includegraphics[width=\linewidth]{CVPR-2026/Figures/cifar_client_gpu_ram_list_comparison.pdf}
    \caption{Client-side RAM usage (MB) over 200 rounds. The chart contrasts the high, stable memory cost of full-model methods ($\sim$600MB) with the low, sub-100MB cost of split-learning methods (SplitFed, HSFL, and Ours).}
    \label{fig:computation}
\end{figure}
\cref{fig:computation} reveals a stark divide: traditional, full-model methods (e.g., FedAvg) require a stable $\sim$600~MB of RAM to load the \textit{entire} model. In contrast, our H-SFP framework and  other split-learning methods consume well under 100 MB. This dramatic reduction is because H-SFP clients only manage the lightweight initial segment ($W_c$). This result confirms ours is highly practical for resource-constrained clients which standard FL is infeasible.

\begin{table*}
\centering
\caption{Ablation study on the components of H-SFP, reporting F1 Score (\%) at client-edge interval $I_c=5$ and edge-cloud interval $I_e=10$. The results show the incremental performance gains from adding contrastive loss (CL) and synthetic data (SD). More importantly, they validate our core design choices by comparing the \textbf{Full H-SFP} against two key ablations: a \textbf{Full ($\mu$-only)} variant (which lacks dimension-aware synthesis) and a \textbf{Full (Purely Generative)} variant (which disables model averaging).}
\begin{tabular}{l|cc|cc|cc}
\hline
\multirow{2}{*}{\textbf{Method}} &
\multicolumn{2}{c|}{\textbf{CIFAR-10}} &
\multicolumn{2}{c|}{\textbf{CIFAR-100}} &
\multicolumn{2}{c}{\textbf{HAM10000}} \\
\cline{2-7}
& \textbf{IID} & $dir(0.7)$ & $\textbf{IID}$ & $dir(0.7)$ & \textbf{IID} & dir(0.7) \\
\hline
Baseline (Split model, no CL) & 52.64\% & 23.20\% & 25.47\% & 0.53\% & 65.12\% & 25.41\% \\
+ Contrastive Loss (CL) & 57.11\% & 37.55\% & 29.73\% & 3.47\% & 69.38\% & 33.58\% \\
+ Synthetic Data (SD) & 60.50\% & 39.62\% & 32.86\% & 4.46\% & 71.57\% & 35.39\% \\
Full ($\mu$-only) & 62.15\% & 41.80\% & 36.40\% & 5.12\% & 74.22\% & 36.10\% \\
Full (Purely Generative, no model avg.) & 55.30\% & 30.15\% & 28.50\% & 2.10\% & 68.40\% & 30.05\% \\
\textbf{Full H-SFP ($\mu$ + $\sigma$ + model avg.)} & \textbf{65.44\%} & \textbf{45.52\%} & \textbf{55.23\%} & \textbf{6.65\%} & \textbf{78.86\%} & \textbf{38.36\%} \\
\hline
\end{tabular}
\label{tab:ablation}
\end{table*}


\subsection{Scalability Test}
    To evaluate robustness to data sparsity, we partition a fixed total dataset across 20 to 200 clients (\cref{tab:main_robustness_sparsity}). As expected, all methods degrade as client count increases, since less data per client leads to significant ``client drift". For instance, FedAvg on CIFAR-10 drops from 72.30\% (20 clients) to 57.25\% (200 clients). Specifically, H-SFP (ours) consistently outperforms all baselines, including the strong FedGen and FedDF methods, at all client counts. This superiority is most evident in the high-sparsity (200 client) scenarios. Here, Ours (A) leads the next-best baseline by up to +7.76\% on HAM10000, +4.94\% on CIFAR-10, and +2.00\% on CIFAR-100. This validates our architecture's resilience. While baselines like FedAvg suffer from aggregating poorly trained, high-variance weights, our generative process uses lightweight feature statistics. This allows edge servers to synthesize a stable cluster-level distribution, effectively filtering noise and mitigating the impact of client drift. Finally, (A) with more frequent aggregation intervals (5,10) performs best, as it corrects for client drift. We note this experiment tests robustness to data sparsity under IID conditions. An analysis of scalability combined with high heterogeneity remains a priority for future investigation.

\subsection{Multi-tier prototype representations}
\label{sec:tsne_analysis}

% --- FIG 1 (TIER 1) ---
\begin{figure}[h]
    \centering
    % Make sure 'tsne_tier1_synthetic.pdf' is in the correct path
    \includegraphics[width=\linewidth]{CVPR-2026/Figures/tsne_tier1_synthetic.pdf}
    \caption{
        T-SNE visualization of the raw feature representations from \textbf{Tier 1 (Client outputs)}. 
        These plots show the combined features from all clients for CIFAR-10, CIFAR-100, and HAM10000 before edge-level aggregation. 
        The representations are heavily overlapped and lack clear class separation, 
        demonstrating the high variance and `client drift' from unaligned local training.
    }
    \label{fig:tsne_tier1}
\end{figure}

% --- FIG 2 (TIER 2) ---
\begin{figure}[h]
    \centering
    % Make sure 'tsne_tier2_synthetic.pdf' is in the correct path
    \includegraphics[width=\linewidth]{CVPR-2026/Figures/tsne_tier2_synthetic.pdf}
    \caption{
        T-SNE visualization of the refined feature representations from \textbf{Tier 2 (Edge outputs)}.
        These plots show the feature space for CIFAR-10, CIFAR-100, and HAM10000 after our proposed generative aggregation at the edge server. 
        The edge server successfully synthesizes and learns from a stable distribution, 
        creating a new feature space with well-separated class clusters.
    }
    \label{fig:tsne_tier2}
\end{figure}

To qualitatively validate the effectiveness of our multi-tier generative process, we visualize the feature representations at different tiers using T-SNE. \cref{fig:tsne_tier1} shows the combined feature outputs from clients across all three datasets. When aggregated from all unaligned clients, they are chaotic and lack any discernible class structure. This offers a challenge of client drift, where each client learns an effective but different local feature space, resulting in an overlapped and nondiscriminative global representation. In contrast, \cref{fig:tsne_tier2} visualizes the feature space after the edge server's generative aggregation. This visual evidence supports our core claim. Instead of propagating error, the generative aggregation acts as a refinement and stabilization mechanism. It successfully filters the chaotic client-level representations (\cref{fig:tsne_tier1}) and transforms them into the highly discriminative, globally consistent feature space shown in \cref{fig:tsne_tier2}.
\subsection{Ablation Studies}
\noindent\textbf{Ablation on H-SFP Components.}
We analyze the contribution of each H-SFP component in \cref{tab:ablation}. The results show clear incremental gains: performance improves step-by-step from the baseline (Split model, no CL) to adding the contrastive loss (+ CL) and further by incorporating synthetic data (+ SD). We then validate two critical design choices of the full model. First, our Full H-SFP model (transmitting $\mu + \sigma$) significantly outperforms the Full ($\mu$-only) variant (e.g., 78.8\% vs. 74.22\% on HAM10000 IID), confirming that transmitting the standard deviation $\sigma$ is essential for our dimension-aware synthesis. Most critically, the Full (Purely Generative) variant (which disables periodic model averaging) suffers a catastrophic performance collapse (e.g., to 2.10\% on CIFAR-100 $dir(0.7)$). This result provides conclusive evidence that our dual-timescale hybrid mechanism is essential for preventing long-term model drift and ensuring stability.

\noindent\textbf{Ablation on Temperature $\tau$.} We ablated the temperature $\tau$ in Appendix B (Tab. 1), finding $\tau=0.5$ provides the best balance of feature discrimination and stability.
% \input{sec/5_comparison_with_HSFL}
\section{Limitations}
    While H-SFP shows strong results, we note three areas for future work. First, our generative synthesis uses a diagonal covariance matrix, which may not capture complex feature correlations; exploring full-covariance models could improve performance. Second, the framework introduces new hyperparameters, such as the temperature $\tau$ and aggregation intervals ($I_c, I_e$), and the optimal selection of these parameters remains an open question. Third, our experimental comparison focuses on architectural (HSFL) and prototype (FedProto) baselines. A direct performance comparison against state-of-the-art non-IID regularization methods, such as SCAFFOLD or MOON, is an important direction for future work. Finally, the privacy guarantees of transmitting generative statistics, while avoiding raw data, require a more formal analysis against potential inference attacks.
\section{Conclusion}
    H-SFP is a three-tier framework decoupling training via a novel generative workflow, achieving state-of-the-art accuracy on non-IID data. Moreover, our experiments validate this ultra-lightweight design, demonstrating a $300\times$ communication reduction over foundational FL baselines and $10 \times$ over strong hybrid approaches. Additionally, hierarchical aggregation stabilizes prototypes, proving it is a stable, robust, and practical FL.
{
    \small
    \bibliographystyle{ieeenat_fullname}
    \bibliography{main}
}

% WARNING: do not forget to delete the supplementary pages from your submission 
% \input{sec/X_suppl}

\end{document}
